%% kapa: exam=1 round=25 year=2014 period=2 form=A subject=appraisal_law subject_ko=부동산관계법규 seq=6 number=6
%% source: Q-Net 감정평가사 기출문제 (cst00309) artlSeq=5101207 file=extracted/1008/5101207/2014년 제25회 감정평가사 1차시험 기출문제/2교시(A형).hwp
%% conversion: kordoc 4.24.51 -> tools/gen1.py (auto); review=proofread-sonnet; flags=none
\begin{question}{6}
국토의 계획 및 이용에 관한 법령상 용도지역에 관하여 옳게 연결된 것은?
\begin{choices}
  \item 제1종전용주거지역 - 저층주택 중심의 양호한 주거환경을 보호하기 위하여 필요한 지역
  \item 제3종일반주거지역 - 중고층주택을 중심으로 편리한 주거환경을 조성하기 위하여 필요한 지역
  \item 준주거지역 - 주거기능을 위주로 이를 지원하는 일부 상업기능 및 공업기능을 보완하기 위하여 필요한 지역
  \item 준공업지역 - 주로 중화학공업, 공해성공업 등을 수용하기 위하여 필요한 지역
  \item 자연녹지지역 - 도시의 자연환경ㆍ경관ㆍ산림 및 녹지공간을 보전할 필요가 있는 지역
\end{choices}
\end{question}
\answer{2}
